# Prisca-GraphRAG: Knowledge Retrieval via Ancient Lineage-Boosted Graph Augmented Generation with Federated Privacy

**Author:** Stephen Paul Lutar Jr.
**Affiliation:** SZL Consulting Ltd
**ORCID:** [0009-0001-0110-4173](https://orcid.org/0009-0001-0110-4173)
**Contact:** stephenlutar2@gmail.com
**Date:** 4 May 2026
**Version:** v5
**Companion papers:** v1 [10.5281/zenodo.19867281](https://doi.org/10.5281/zenodo.19867281), v2 [10.5281/zenodo.19944926](https://doi.org/10.5281/zenodo.19944926), v3 (Lutar Invariant), v4 (Lutar Omega Formalism)
**Reference implementation:** [github.com/szl-holdings/szl-holdings-platform](https://github.com/szl-holdings/szl-holdings-platform) -- `packages/ouroboros-integrations/src/sovereign-engine.ts` (Prisca-GraphRAG, VOTE-RAG, FPP classes)
**License:** CC BY 4.0

---

## Abstract

We introduce Prisca-GraphRAG, a retrieval-augmented generation architecture that structures knowledge graphs according to *prisca theologia* --- the Renaissance doctrine of a single ancient theological tradition flowing through Hermes Trismegistus, Zoroaster, Orpheus, Pythagoras, and Plato. Each knowledge node carries a lineage vector encoding its epistemic provenance, and retrieval scoring combines semantic similarity with lineage coherence weighted by the Lutar Omega signature (v4). We extend this with VOTE-RAG, a multi-retriever ensemble that aggregates ranked lists via Borda count weighted by Omega scores. A Federated Prisca Privacy (FPP) layer enables cross-organizational knowledge sharing without exposing raw embeddings, using differential-privacy noise calibrated to the Bekenstein information bound.

This paper does **not** claim empirical superiority over standard RAG benchmarks. The contribution is the formal architecture and its public reference implementation.

---

## 1. Motivation and Continuity from v3--v4

The v3 Lutar Invariant established axiomatic trust aggregation; v4 extended this to energy-mass-information signatures. The present paper addresses a complementary problem: how to structure the *knowledge retrieval* pipeline so that the provenance of retrieved information is formally tracked and scored.

Standard RAG (Lewis et al., 2020) treats all documents as epistemically equivalent. GraphRAG (Edge et al., 2024) adds structural relationships but not epistemic provenance. We propose that the *prisca theologia* tradition provides a principled framework for organizing knowledge by transmission lineage.

---

## 2. Prisca-GraphRAG Architecture

### 2.1 Prisca Lineage Graph

A Prisca Lineage Graph \( G = (V, E, \lambda) \) consists of:

- **Nodes** \( V \): knowledge chunks with embedding vectors \( \mathbf{e}_v \in \mathbb{R}^d \)
- **Edges** \( E \): directed provenance relationships
- **Lineage function** \( \lambda: V \to \{0,1\}^{|\mathcal{L}|} \) mapping each node to lineage membership

The six canonical lineages: Hermetic (Egyptian), Zoroastrian (Persian), Orphic (Greek), Pythagorean (Mathematical), Platonic (Philosophical), Newtonian (Modern-Scientific).

### 2.2 Retrieval Scoring

Given a query \( q \) with embedding \( \mathbf{e}_q \) and inferred lineage vector \( \lambda_q \):

\[
\text{score}(q, v) = \alpha \cdot \text{sim}(\mathbf{e}_q, \mathbf{e}_v) + (1-\alpha) \cdot \text{overlap}(\lambda_q, \lambda_v) \cdot L_\Omega(v)
\]

where \( \text{sim} \) is cosine similarity, \( \text{overlap} \) is the Jaccard coefficient of lineage vectors, \( L_\Omega(v) \) is the v4 Lutar Omega quality signature of the chunk, and \( \alpha \in [0,1] \) is a mixing parameter.

---

## 3. VOTE-RAG Ensemble

Given \( k \) retrievers each returning a ranked list of \( n \) candidates, the VOTE-RAG score for candidate \( v \) is:

\[
\text{VOTE}(v) = \sum_{i=1}^{k} w_i \cdot (n - \text{rank}_i(v) + 1)
\]

where \( w_i = L_\Omega(\text{retriever}_i) \) weights each retriever by its Omega quality score. This combines the Borda count aggregation mechanism from social choice theory with the v4 Omega quality metric.

---

## 4. Federated Prisca Privacy (FPP)

The FPP protocol enables multi-party knowledge sharing while preserving embedding privacy:

1. Each participant computes local prisca lineage vectors and Omega signatures.
2. Lineage vectors are shared in cleartext (categorical, low-dimensional).
3. Embedding similarities are computed via secure inner product protocols.
4. Differential privacy noise \( \eta \sim \mathcal{N}(0, \sigma^2) \) is added with \( \sigma \) calibrated to the Bekenstein bound (v4, Section 3).
5. Aggregate retrieval scores are computed by a coordinator without accessing raw embeddings.

**Privacy guarantee:** Under \( (\epsilon, \delta) \)-differential privacy with \( \sigma = \sqrt{2 \ln(1.25/\delta)} / \epsilon \), the FPP protocol satisfies:

\[
\Pr[\mathcal{M}(D) \in S] \leq e^\epsilon \cdot \Pr[\mathcal{M}(D') \in S] + \delta
\]

for any adjacent datasets \( D, D' \) differing in one participant's embeddings.

---

## 5. Bekenstein-Gated Retrieval

Before returning retrieved chunks to the generator, each candidate passes through the Bekenstein Gate (v4, Section 3):

\[
\text{pass}(v) = \mathbb{1}\left[ I(v) \leq \frac{2\pi R_v E_v}{\hbar c \ln 2} \right]
\]

Chunks exceeding their information-theoretic capacity relative to their provenance chain length are flagged as potentially fabricated.

---

## 6. Connection to v3 Lutar Invariant

The retrieval pipeline feeds into the v3 Lutar Invariant as follows:

- **Axis 1 (data freshness):** determined by the timestamp of retrieved chunks
- **Axis 2 (source priority):** determined by the lineage coherence score
- **Axis 3 (validator pass):** determined by the Bekenstein Gate result

The remaining six axes are computed downstream by the trust evaluation pipeline. This ensures end-to-end traceability from retrieval through aggregation to trust scoring.

---

## 7. What this paper does and does not claim

**Claims:**
1. The Prisca lineage structure is well-defined and maps cleanly to historical transmission chains documented in Renaissance scholarship (Yates, 1964).
2. The VOTE-RAG scoring function is a valid Borda aggregation with Omega weights.
3. The FPP protocol satisfies standard differential privacy definitions.
4. All components are implemented in the public reference implementation.

**Non-claims:**
- No claim of empirical superiority over standard RAG on benchmarks like NaturalQuestions or TriviaQA.
- No claim that lineage labels are objective truth; they are a taxonomic framework.
- No claim of deployed production use.

---

## 8. Conclusion

Prisca-GraphRAG demonstrates that epistemically structured knowledge graphs, organized according to historical transmission lineages, provide a principled framework for retrieval scoring. The VOTE-RAG ensemble and Federated Prisca Privacy extensions provide practical scalability. The framework integrates with the v3 Lutar Invariant and v4 Omega signatures to maintain end-to-end trust traceability.

---

## References

1. Lewis, P. et al. (2020). Retrieval-augmented generation for knowledge-intensive NLP tasks. *NeurIPS*.
2. Edge, D. et al. (2024). From local to global: A graph RAG approach. *arXiv:2404.16130*.
3. Ficino, M. (1474). *De Christiana Religione*.
4. Yates, F.A. (1964). *Giordano Bruno and the Hermetic Tradition*. University of Chicago Press.
5. Dwork, C. et al. (2006). Calibrating noise to sensitivity in private data analysis. *TCC*.
6. Lutar, S.P. (2026a). *The Ouroboros Thesis* (v1). Zenodo. [DOI 10.5281/zenodo.19867281](https://doi.org/10.5281/zenodo.19867281).
7. Lutar, S.P. (2026b). *The Loop Is the Product* (v2). Zenodo. [DOI 10.5281/zenodo.19944926](https://doi.org/10.5281/zenodo.19944926).
8. Lutar, S.P. (2026c). *The Lutar Invariant* (v3). [github.com/szl-holdings/ouroboros-thesis](https://github.com/szl-holdings/ouroboros-thesis).
9. Lutar, S.P. (2026d). *The Lutar Omega Formalism* (v4). This series.
